% !TeX program = xelatex
% Bien dich bang XeLaTeX (Overleaf: Menu > Compiler > XeLaTeX).
% Diem ky nang va vai tro theo noi dung nguoi viet da hoan thien.
% So lieu Housing: docs/integration/quangdung-housing-performance-handoff-2026-10-10.md
% Phan chu nha: apps/web/src/landlord/; apps/platform-api/src/housing-management.ts;
% housing-revenue.ts; docs/review/landlord-revenue-admin-glass-2026-10-17.md.
% Doanh thu la UOC TINH theo ky thue, khong phai so tien da doi soat ngan hang.
% Moc truoc/sau va su co cache theo thong tin nguoi viet bo sung.
\documentclass[10pt,a4paper]{article}
\usepackage[margin=13mm,top=12mm,bottom=8mm]{geometry}
\usepackage{fontspec}
\IfFontExistsTF{Arial}{\setmainfont{Arial}}{\setmainfont{TeX Gyre Heros}}
\usepackage{xcolor}
\definecolor{ink}{HTML}{16324F}
\definecolor{accent}{HTML}{137C83}
\definecolor{muted}{HTML}{556270}
% Lien ket PDF truc tiep cua XeLaTeX, khong can cai them hyperref.
\newcommand{\href}[2]{\leavevmode\special{pdf:bann << /Type /Annot /Subtype /Link /Border [0 0 0] /A << /S /URI /URI (#1) >> >>}{\color{accent}#2}\special{pdf:eann}}
\pagestyle{empty}
\setlength{\parindent}{0pt}
\setlength{\parskip}{2pt}
\setlength{\emergencystretch}{2em}
\newcommand{\can}[1]{\textcolor{accent}{\textit{[#1]}}}
\newcommand{\cauHoi}[2]{\par\addvspace{3pt}{\bfseries\color{ink}#1. #2}\par\nopagebreak}
\newenvironment{cacY}{\begin{list}{\textcolor{ink}{\textbullet}}{
  \setlength{\leftmargin}{12pt}\setlength{\labelsep}{5pt}
  \setlength{\itemsep}{1pt}\setlength{\parsep}{0pt}\setlength{\topsep}{2pt}
}}{\end{list}}

\begin{document}
\fontsize{10.2}{12}\selectfont

{\centering
{\fontsize{14}{17}\selectfont\bfseries\color{ink} Giới thiệu bản thân sau 6 tuần học chương trình AI thực chiến }
\par}
\vspace{3pt}

Em là \textbf{Mai Quang Dũng}, mã học viên \textbf{2A202602966}, sinh viên ngành Khoa học dữ liệu,
Trường ĐH Khoa học Tự nhiên, ĐHQG-HCM. Em định hướng phát triển hệ thống AI
\textbf{end-to-end}, từ dữ liệu, mô hình/agent đến backend, giao diện và vận hành;
muốn học sâu DevOps/MLOps để theo sát vòng đời sản phẩm.

\textbf{Repository:} \href{https://github.com/Dungkdl/ReloHome}{github.com/Dungkdl/ReloHome}

\cauHoi{1}{Trước và sau 6 tuần, em làm thêm được gì?}
\begin{cacY}
\item \textbf{Trước khóa học:} Em đã làm RAG (sử dụng hybrid retrieval, reranking, v.v.)
và bộ đánh giá 50 ca; có kinh nghiệm SQL, fine-tuning/domain adaptation và
distillation cho FinBERT trong dự án tài chính FinSight.
\item \textbf{Sau khóa học:} Em có thêm kinh nghiệm thiết kế agent điều phối đa dịch vụ
và MCP. Ví dụ: Tìm nhà qua \texttt{search\_housing}, chuyển địa chỉ đã chọn sang báo giá
chuyển đồ; các tác vụ độc lập chạy song song, tác vụ phụ thuộc chờ kết quả trước.
\end{cacY}

\cauHoi{2}{Sản phẩm tốt nhất trong khóa: ReloHome -- Top 10 Demo khóa 4}
\begin{cacY}
\item \textbf{Bài toán và đóng góp:} ReloHome kết nối các nhu cầu chuyển nơi ở.
Em thiết kế dịch vụ chuyển nhà, ghép trọ và CRM chủ nhà: quản lý nhà/phòng,
tin đăng, khách liên hệ, tin nhắn, lịch xem, báo chuyển tiền và doanh thu ước tính.
\item \textbf{Kiến trúc và đo lường:} Em hỗ trợ thiết kế agent với LangGraph/MCP,
tối ưu thực thi tuần tự/song song; dùng Langfuse theo dõi trace, độ trễ và token
để phân tích lỗi, chi phí và hỗ trợ đánh giá chất lượng.
\item \textbf{Cách dùng AI:} Em dùng AI phác thảo sơ đồ kiến trúc ban đầu, rồi dựa
vào đó thiết kế chi tiết ranh giới dịch vụ, luồng dữ liệu; điều chỉnh phương án
qua đối chiếu yêu cầu, mã nguồn và kiểm thử.
\item \textbf{Kết quả:} Thời gian SQL tìm kiếm Housing giảm từ \textbf{274,18 xuống 28,33 ms}
trong benchmark local; kết quả được đối chiếu trước và sau tối ưu.
\item \textbf{Nếu làm lại:} Em sẽ xây benchmark sớm cho chất lượng câu trả lời, độ đúng
của tool và độ trễ; cho người dùng trải nghiệm lâu hơn để thu thập phản hồi và lỗi thực tế.
\end{cacY}

\cauHoi{3}{Một lần AI làm sai và bài học rút ra}
\begin{cacY}
\item AI tăng tốc truy vấn bằng cache, chạy local nhanh hơn nhưng gây lỗi khi triển khai production.
\item Em phát hiện hệ thống phản hồi rất chậm trên production dù chạy nhanh ở local.
Em xử lý bằng cách thiết kế lại các node trong LangGraph và kiểm tra kỹ kiến trúc,
thay vì dựa hoàn toàn vào AI.
\item Bài học: tối ưu phải giữ tính đúng đắn thay vì chỉ nhìn thời gian query.
\end{cacY}

\cauHoi{4}{Điều khó nhất hoặc em chưa hiểu rõ}
\begin{cacY}
\item Khó nhất là khoảng cách giữa kiến trúc được học và thực tế: sơ đồ hợp lý vẫn
phải xử lý dữ liệu thay đổi, quyền truy cập, timeout và trạng thái liên dịch vụ.
Chạy đúng ở local chưa bảo đảm chạy đúng trên production.
\item Em từng bước đối chiếu contract, log và kiểm thử tích hợp để tìm khác biệt.
Em cần học thêm cấu hình môi trường, CI/CD, giám sát và kiểm thử gần điều kiện vận hành thật.
\end{cacY}

\cauHoi{5}{Cách em làm việc nhóm và xử lý bất đồng}
\begin{cacY}
\item Em thường phụ trách thiết kế, xây dựng hệ thống, hỗ trợ kiểm thử và nghiên cứu tài liệu.
\item Nhóm từng khó thống nhất thiết kế frontend vì chưa ai có nền tảng; cả nhóm
họp, trao đổi kỹ các phương án để thống nhất hướng triển khai.
\end{cacY}

\cauHoi{6}{Vai trò em muốn nhận, theo thứ tự ưu tiên}
\begin{cacY}
\item \textbf{Ưu tiên 1 -- Agent (lập trình agent, MCP):} Đã thiết kế agent và tích hợp MCP;
em muốn học sâu về quản lý trạng thái, đánh giá và xử lý lỗi.
\item \textbf{Ưu tiên 2 -- Ngữ cảnh \& test (rule cho AI, test E2E):} Đã kiểm thử backend/API
và luồng người dùng; em muốn học sâu về rule cho AI và E2E.
\item Nếu có cơ hội, em mong được thử sức thêm với \textbf{Hạ tầng (Docker, GitLab, DB)
và CI/CD} để hiểu rõ hơn việc triển khai, tự động hóa và vận hành hệ thống. Đây cũng là phần mà bản thân em qua dự án vừa rồi cảm thấy khó khăn nhất.
\end{cacY}

\newpage
\textbf{Tự đánh giá kỹ năng}
{\small\color{muted}(Thang điểm 0--3.)}
\par\vspace{2pt}
{\fontsize{9.2}{10.5}\selectfont
\setlength{\tabcolsep}{4pt}
\renewcommand{\arraystretch}{1.08}
\begin{tabular}{@{}p{34mm}p{11mm}p{\dimexpr\linewidth-45mm-4\tabcolsep\relax}@{}}
\hline
\textbf{\color{ink}Kỹ năng} & \textbf{\color{ink}Điểm} & \textbf{\color{ink}Minh chứng} \\
\hline
Git & \multicolumn{1}{c}{2/3} & Làm việc trên nhánh riêng; biết push, rebase, merge và tạo pull request. \\
Linux & \multicolumn{1}{c}{0/3} & Đã biết tới nhưng chưa từng dùng. \\
Docker & \multicolumn{1}{c}{2/3} & Container hóa ứng dụng RAG và đóng gói Location MCP. \\
SQL & \multicolumn{1}{c}{2/3} & Xử lý dữ liệu PostGIS, tối ưu Housing bằng cách đánh index có đo và đối chiếu kết quả. \\
Java & \multicolumn{1}{c}{0/3} & Đã biết tới nhưng chưa từng dùng. \\
TypeScript & \multicolumn{1}{c}{2/3} & Đóng góp dịch vụ nhà ở, MCP và giao diện chủ nhà. \\
Python & \multicolumn{1}{c}{2/3} & Làm pipeline ML, RAG và backend/agent Roommate. \\
Prompt/Codex & \multicolumn{1}{c}{2/3} & Sử dụng thành thạo Codex ở mức thực hành; biết dùng thêm plugin và skill. \\
Agent/MCP & \multicolumn{1}{c}{2/3} & Làm Agentic RAG, Location MCP và tích hợp Roommate vào agent. \\
Viết test & \multicolumn{1}{c}{2/3} & Đóng góp kiểm thử backend/API và Roommate mobile (\texttt{92f04cb}). \\
CI/CD & \multicolumn{1}{c}{1/3} & Sửa phạm vi dữ liệu trong image migration (\texttt{434f5a8}). \\
Viết tài liệu & \multicolumn{1}{c}{2/3} & Viết hướng dẫn dữ liệu, bàn giao và ghi nhận kết quả kiểm thử. \\
\hline
\end{tabular}}

Em cảm ơn anh đã xem hết ạ.

\end{document}
